\documentclass[12pt]{article}
\usepackage[margin=1in]{geometry}
\usepackage{amsmath,amssymb}
\usepackage{graphicx}
\usepackage{booktabs}
\usepackage{hyperref}
\usepackage{natbib}
\usepackage{lineno}
\usepackage{xcolor}

\title{Profiling of Chimeric RNAs in Human Retinal Development with Retinal Organoids}
\author{Author One$^{1}$, Author Two$^{1}$, Author Three$^{2}$, Author Four$^{1,*}$\\
\\
\small $^{1}$Department of Ophthalmology, University Medical Center\\
\small $^{2}$Institute of Developmental Biology\\
\small $^{*}$Corresponding author}
\date{}

\begin{document}
\maketitle

\begin{abstract}
Chimeric RNAs have emerged as important regulators of stem and progenitor cell differentiation in central nervous system development, yet their landscape in human retinal development remains completely uncharacterized. Here, we present the first comprehensive atlas of chimeric RNAs across human retinal organoid (RO) development, profiling 22 samples spanning eight developmental stages from day 0 (D0) to day 120 (D120). Using RNA sequencing and FusionCatcher-based chimeric RNA detection, we found that intra-chromosomal chimeric RNAs increase significantly with developmental progression ($r^2 = 0.93$, $p = 0.00076$, Pearson correlation), while inter-chromosomal events do not follow this trend. The majority of chimeric RNAs (61.9\%) were formed between two protein-coding genes, with in-frame fusions accounting for 11.2\% of all events. We identified and validated four isoforms of the chimeric RNA CTNNBIP1-CLSTN1 (CTCL)---previously implicated in cerebral organoid development---in retinal organoids via PCR and Sanger sequencing. Lentiviral shRNA-mediated knockdown of CTCL obstructed retinal organoid neural differentiation but unexpectedly promoted retinal pigment epithelium (RPE) fate, with transcriptomic analysis revealing alterations in epithelial-mesenchymal transition and extracellular matrix pathways rather than the anticipated Wnt signaling changes. These findings establish the chimeric RNA transcriptome of the developing human retina and reveal a novel role for CTCL in balancing neural retinal and RPE cell fate decisions during retinogenesis.
\end{abstract}

\section{Introduction}

Chimeric RNAs---transcripts composed of sequences from two or more distinct genomic loci---have been increasingly recognized as functional molecules in both normal development and disease \citep{li2009chimeric,wu2019chimeric}. Originally identified primarily in cancer contexts as products of chromosomal translocations, chimeric RNAs are now known to arise through various mechanisms including cis-splicing, trans-splicing, and read-through transcription in non-pathological cells \citep{kannan2011recurrent,jividen2014chimeric}. In the central nervous system (CNS), chimeric RNAs have been shown to regulate stem cell maintenance and neuronal differentiation, suggesting they constitute an underappreciated layer of gene regulation during neural development \citep{wu2019chimeric,chase2020role}.

The retina, as a specialized extension of the CNS, undergoes a highly orchestrated developmental program in which retinal progenitor cells give rise to six neuronal and one glial cell type in a conserved temporal order \citep{cepko2014intrinsically,livesey2001vertebrate}. Understanding the molecular mechanisms governing retinal cell fate specification is essential for developing therapeutic strategies for retinal degenerative diseases and for optimizing retinal organoid differentiation protocols for regenerative medicine \citep{eiraku2011self,nakano2012self}. However, while transcriptomic and epigenomic landscapes of retinal development have been extensively profiled \citep{clark2019single,lu2020single}, the chimeric RNA transcriptome of the developing retina has never been catalogued.

Human embryonic stem cell (hESC)-derived retinal organoids (ROs) provide a powerful model system for studying human retinal development, recapitulating the temporal sequence of retinal cell type generation from progenitors through photoreceptor maturation over a 120-day culture period \citep{zhong2014generation,capowski2019reproducibility}. ROs overcome the practical and ethical limitations of obtaining staged human retinal tissue, enabling systematic molecular profiling across developmental stages.

We previously demonstrated that the chimeric RNA CTNNBIP1-CLSTN1 (CTCL) plays a critical role in human cerebral organoid development, where its downregulation leads to premature neuronal differentiation \citep{wu2019chimeric}. CTNNBIP1 encodes an inhibitor of the Wnt/$\beta$-catenin signaling pathway, while CLSTN1 is involved in vesicular transport. Whether CTCL or other chimeric RNAs function during retinogenesis has not been investigated.

In this study, we address three questions: (1) What is the global landscape of chimeric RNAs during human retinal organoid development? (2) Is the chimeric RNA CTCL expressed in developing retinal organoids, and if so, in what isoforms? (3) What is the functional consequence of CTCL loss-of-function on retinal organoid differentiation? Our findings reveal a dynamic chimeric RNA landscape across retinal development and an unexpected role for CTCL in directing the balance between neural retinal and RPE cell fates.

\section{Methods}

\subsection{Retinal Organoid Generation and Culture}

Human embryonic stem cells (hESCs) carrying a CRX-tdTomato reporter were maintained on Matrigel-coated plates in mTeSR1 medium and differentiated into retinal organoids following established protocols \citep{zhong2014generation}. Briefly, hESCs were dissociated into small clumps and cultured in suspension to form embryoid bodies, followed by adherent culture with neural induction medium. Optic vesicle-like structures were manually excised and cultured in suspension in retinal differentiation medium supplemented with taurine and retinoic acid. Organoids were maintained through day 120 (D120) with medium changes every 2--3 days. Samples were collected at eight developmental stages: D0, D30, D45, D60, D75, D90, D105, and D120, with 2--3 biological replicates per stage totaling 22 samples.

\subsection{RNA Sequencing and Chimeric RNA Detection}

Total RNA was extracted using TRIzol reagent and quality assessed by Agilent Bioanalyzer (RIN $>$ 7.0 for all samples). Libraries were prepared using the VAHTS Universal V6 RNA-seq Library Prep Kit for Illumina and sequenced on an Illumina NovaSeq platform with paired-end 150 bp reads. Chimeric RNAs were identified using FusionCatcher software \citep{nicorici2014fusioncatcher} with a stringency filter requiring at least one spanning unique read. Expression levels were quantified as $\log(\text{spanning unique reads} / \text{total reads} \times 10^7)$. For parental gene expression analysis, reads were mapped to the hg38 reference genome using Hisat2 \citep{kim2019graph} and transcript abundances were determined with featureCounts \citep{liao2014featurecounts}.

\subsection{Chimeric RNA Characterization}

Chimeric RNAs were classified by genomic distribution (intra-chromosomal versus inter-chromosomal), predicted effect on the reading frame (in-frame, CDS truncated, UTR-UTR, etc.), and parental gene biotype (protein-coding, lncRNA, pseudogene, etc.). Sequence motifs around fusion sites were analyzed using 20 bp flanking sequences with the seqLogo R package. Correlation between chimeric RNA expression and parental gene expression was assessed by Spearman correlation.

\subsection{CTCL Isoform Validation}

CTCL isoforms were validated by RT-PCR using primers spanning the fusion junction, followed by Sanger sequencing. Four distinct isoforms were amplified from D60 retinal organoid cDNA: an in-frame isoform (CTNNBIP1 exons 1--5 fused to CLSTN1 terminal exons), two UTR-CDS truncated variants, and one CDS-intronic variant.

\subsection{CTCL Knockdown Experiments}

Lentiviral particles encoding shRNA targeting the CTCL fusion junction (shCTCL) or scramble control were produced in HEK293T cells. ROs were transduced at D12 (multiplicity of infection $\approx$ 5), and knockdown efficiency was confirmed by RT-qPCR at 48--72 hours post-infection. Three independent experiments were performed. Organoids were collected at D60 for downstream analyses.

\subsection{Immunohistochemistry}

D60 organoids were fixed in 4\% paraformaldehyde, cryoprotected in 30\% sucrose, embedded in OCT, and sectioned at 14 $\mu$m. Sections were stained with antibodies against CRX, OTX2, VSX2, RCVRN, MITF, RPE65, and PAX6 following standard immunofluorescence protocols. Images were acquired on a Zeiss LSM 800 confocal microscope.

\subsection{Transcriptomic Analysis of Knockdown Organoids}

RNA from shCTCL and scramble D60 organoids was sequenced as above. Differentially expressed genes (DEGs) were identified using criteria of $|\log_2(\text{shCTCL/scramble})| > 1$ and $p < 0.05$. Functional enrichment of DEGs was performed using Metascape \citep{zhou2019metascape}. Gene Set Enrichment Analysis (GSEA) was conducted using the R package clusterProfiler \citep{yu2012clusterprofiler} with the c8 cell type signature reference gene set. Parental gene expression (CTNNBIP1 and CLSTN1) was assessed by RT-qPCR in knockdown versus control organoids ($n = 4$ per group).

\section{Results}

\subsection{Chimeric RNA Landscape Across Retinal Organoid Development}

To characterize the chimeric RNA transcriptome of the developing human retina, we performed RNA sequencing on 22 retinal organoid samples spanning eight developmental stages from D0 to D120 (Table~\ref{tab:samples}). Using FusionCatcher with a minimum threshold of one spanning unique read, we detected chimeric RNA events at every developmental stage, establishing that chimeric RNAs are continuously expressed throughout retinogenesis.

\begin{table}[htbp]
\centering
\caption{Retinal organoid sample characteristics and chimeric RNA detection parameters.}
\label{tab:samples}
\begin{tabular}{lll}
\toprule
Parameter & Value \\
\midrule
Cell line & hESC-derived (CRX-tdTomato reporter) \\
Total RO samples sequenced & 22 \\
Developmental stages & 8 (D0, D30, D45, D60, D75, D90, D105, D120) \\
Chimeric RNA detection tool & FusionCatcher \\
Positive criteria & Spanning unique reads $\geq$ 1 \\
Expression metric & $\log(\text{spanning unique reads}/\text{total reads} \times 10^7)$ \\
Genome reference & hg38 \\
Alignment tool & Hisat2 \\
Transcript assembly & featureCounts \\
\bottomrule
\end{tabular}
\end{table}

Analysis of chimeric RNA genomic distribution revealed a striking divergence between intra-chromosomal and inter-chromosomal events. Intra-chromosomal chimeric RNAs increased significantly with developmental progression, showing a strong positive correlation with organoid age ($r^2 = 0.93$, $p = 0.00076$, Pearson correlation). In contrast, inter-chromosomal chimeric RNAs showed no significant developmental trend (Table~\ref{tab:trends}). This observation suggests that the mechanisms generating intra-chromosomal chimeric RNAs---likely cis-splicing or read-through transcription---become more active or permissive as retinal differentiation proceeds, potentially reflecting increased transcriptional complexity at later developmental stages.

\begin{table}[htbp]
\centering
\caption{Developmental trends of chimeric RNA categories in retinal organoids.}
\label{tab:trends}
\begin{tabular}{lll}
\toprule
Chimeric RNA Category & Trend Over Development & Statistical Test \\
\midrule
Intra-chromosomal & Increase with development & $r^2 = 0.93$, $p = 0.00076$ (Pearson) \\
Inter-chromosomal & No significant trend & Not significant \\
Total chimeric events & Continuously expressed & Present at all stages \\
\bottomrule
\end{tabular}
\end{table}

\subsection{Characterization of Chimeric RNA Properties}

We next examined the predicted functional consequences of the detected chimeric RNAs. Classification by predicted effect on the reading frame revealed that in-frame chimeric RNAs---those predicted to produce fusion proteins---accounted for 11.2\% of all events (Table~\ref{tab:characterization}). CDS(truncated)-UTR and UTR-UTR fusions each represented approximately 10.7\% of events, with the remaining chimeric RNAs distributed across 11 additional categories.

\begin{table}[htbp]
\centering
\caption{Characterization of chimeric RNAs by predicted effect and parental gene biotype.}
\label{tab:characterization}
\begin{tabular}{ll}
\toprule
Predicted Effect Category & Percentage \\
\midrule
In frame & 11.2\% \\
CDS(truncated)\_UTR & 10.7\% \\
UTR\_UTR & 10.7\% \\
Other categories (11 more) & Remaining \\
\midrule
\textbf{Parental Gene Biotype Combination} & \textbf{Percentage} \\
\midrule
Protein-coding + Protein-coding & 61.9\% \\
Other combinations & 38.1\% \\
\bottomrule
\end{tabular}
\end{table}

Analysis of parental gene biotypes showed that 61.9\% of chimeric RNAs were formed between two protein-coding genes, consistent with the predominance of protein-coding transcription in the developing retina. Sequence motif analysis of 20 bp flanking the fusion sites did not reveal strongly conserved consensus sequences, suggesting that chimeric RNA generation in retinal organoids is not driven by specific sequence motifs. Importantly, Spearman correlation analysis showed no significant relationship between chimeric RNA expression levels and parental gene expression ($p > 0.05$), indicating that chimeric RNA abundance is not simply a byproduct of high parental gene transcription but is regulated independently.

\subsection{CTCL Isoform Identification in Retinal Organoids}

Given our prior finding that the chimeric RNA CTCL regulates cerebral organoid development \citep{wu2019chimeric}, we investigated whether CTCL is expressed in retinal organoids. RNA-seq analysis detected CTCL spanning unique reads across multiple developmental stages. RT-PCR using junction-spanning primers followed by Sanger sequencing confirmed four distinct CTCL isoforms in D60 retinal organoids (Table~\ref{tab:isoforms}): (1) an in-frame isoform in which CTNNBIP1 exons 1--5 are fused to the last 17 exons of CLSTN1, predicted to encode a chimeric protein; (2--3) two UTR-CDS(truncated) variants with truncated coding regions; and (4) a CDS(truncated)-intronic variant incorporating intronic sequence from the CLSTN1 locus.

\begin{table}[htbp]
\centering
\caption{CTCL isoforms identified in retinal organoids.}
\label{tab:isoforms}
\begin{tabular}{llp{7cm}}
\toprule
Isoform & Type & Description \\
\midrule
1 & In-frame & CTNNBIP1 exons 1--5 fused to CLSTN1 (last 17 exons) \\
2 & UTR\_CDS(truncated)-1 & Truncated coding region \\
3 & CDS(truncated)\_Intronic & Intronic fusion variant \\
4 & UTR\_CDS(truncated)-2 & Second truncated variant \\
\bottomrule
\end{tabular}
\end{table}

Heatmap analysis of CTCL spanning unique reads across developmental stages revealed dynamic expression patterns, with the in-frame isoform showing the highest expression at intermediate stages (D45--D75), coinciding with the period of active retinal progenitor cell differentiation.

\subsection{CTCL Knockdown Obstructs Retinal Differentiation but Promotes RPE Fate}

To assess CTCL function in retinogenesis, we performed lentiviral shRNA-mediated knockdown at D12 and analyzed organoids at D60 (Table~\ref{tab:knockdown_design}). Knockdown efficiency was confirmed at 48--72 hours post-infection. Importantly, RT-qPCR analysis showed that shCTCL did not affect the expression of either parental gene (CTNNBIP1, $p > 0.05$; CLSTN1, $p > 0.05$; $n = 4$), confirming specificity of the knockdown for the chimeric transcript.

\begin{table}[htbp]
\centering
\caption{CTCL knockdown experimental design.}
\label{tab:knockdown_design}
\begin{tabular}{ll}
\toprule
Parameter & Value \\
\midrule
Reporter line & CRX-tdTomato \\
shRNA delivery & Lentiviral transduction at D12 \\
Controls & Scramble shRNA \\
Independent experiments & 3 \\
Collection timepoint & D60 \\
Knockdown efficiency assessment & 48--72 hours post-infection \\
\bottomrule
\end{tabular}
\end{table}

Morphological analysis revealed that shCTCL organoids exhibited thinner outer layers compared to scramble controls at D60, suggesting impaired lamination. Immunohistochemistry revealed a complex phenotype (Table~\ref{tab:ihc}): neural retinal markers CRX and RCVRN (photoreceptor precursors) were reduced in shCTCL organoids, while RPE-specific markers MITF and RPE65 were markedly increased. The retinal progenitor marker VSX2 showed altered distribution, and PAX6 staining revealed an abnormal pattern. These results indicate that CTCL knockdown simultaneously obstructs neural retinal differentiation and promotes RPE cell fate---a striking bifurcation in developmental trajectory.

\begin{table}[htbp]
\centering
\caption{Immunohistochemistry results in CTCL-knockdown organoids at D60.}
\label{tab:ihc}
\begin{tabular}{lll}
\toprule
Marker & Category & shCTCL vs Control \\
\midrule
CRX & Photoreceptor precursor & Reduced \\
OTX2 & Neural/retinal marker & Altered \\
VSX2 & Retinal progenitor & Changed \\
RCVRN & Photoreceptor & Reduced \\
MITF & RPE-specific & Increased \\
RPE65 & RPE-specific & Increased \\
PAX6 & Eye field/progenitor & Altered pattern \\
\bottomrule
\end{tabular}
\end{table}

\subsection{Transcriptomic Analysis Reveals Unexpected Pathway Alterations}

To elucidate the molecular mechanisms underlying the CTCL knockdown phenotype, we performed RNA-seq on shCTCL and scramble D60 organoids. Differential expression analysis identified hundreds of DEGs ($|\log_2\text{FC}| > 1$, $p < 0.05$). Functional enrichment of upregulated genes revealed significant enrichment for epithelial-mesenchymal transition (EMT), extracellular matrix organization, and epithelial cell differentiation pathways. Downregulated genes were enriched for neurogenesis and morphogenesis-related terms.

\begin{table}[htbp]
\centering
\caption{Expected versus observed pathway changes in CTCL-knockdown organoids.}
\label{tab:pathways}
\begin{tabular}{lll}
\toprule
Pathway & Expected? & Observed? \\
\midrule
Wnt signaling modulation & Yes (CTNNBIP1 function) & No significant change \\
Vesicular transport & Yes (CLSTN1 function) & No significant change \\
EMT pathway & Not expected & Significantly changed \\
Extracellular matrix organization & Not expected & Significantly changed \\
Epithelial cell differentiation & Not expected & Significantly changed \\
\bottomrule
\end{tabular}
\end{table}

Strikingly, neither Wnt signaling (expected based on the CTNNBIP1 parental gene function) nor vesicular transport pathways (expected based on CLSTN1 function) showed significant modulation (Table~\ref{tab:pathways}). This dissociation between parental gene function and chimeric RNA function provides strong evidence that CTCL acts as a functionally distinct entity with emergent biological activities not predictable from its constituent genes.

GSEA using the c8 cell type signature reference gene set revealed significant enrichment for RPE-related gene signatures in shCTCL organoids. Consistent with the immunohistochemistry findings, RPE marker genes including MITF, RPE65, BEST1, and TYR were all significantly upregulated in knockdown organoids (Table~\ref{tab:rpe}).

\begin{table}[htbp]
\centering
\caption{RPE marker gene expression in CTCL-knockdown organoids.}
\label{tab:rpe}
\begin{tabular}{lll}
\toprule
RPE Marker Gene & Direction in shCTCL & Significance \\
\midrule
MITF & Upregulated & Significant \\
RPE65 & Upregulated & Significant \\
BEST1 & Upregulated & Significant \\
TYR & Upregulated & Significant \\
\bottomrule
\end{tabular}
\end{table}

\section{Discussion}

This study presents the first comprehensive characterization of chimeric RNAs during human retinal development using a retinal organoid model system. Our findings reveal several novel aspects of chimeric RNA biology in the developing retina that have implications for understanding retinal cell fate decisions and the broader role of chimeric RNAs in CNS development.

The observation that intra-chromosomal chimeric RNAs increase with developmental progression ($r^2 = 0.93$, $p = 0.00076$) while inter-chromosomal events remain stable suggests mechanistic differences in chimeric RNA biogenesis. Intra-chromosomal chimeric RNAs likely arise from cis-splicing or read-through transcription events that may be facilitated by increasing transcriptional activity and chromatin remodeling during differentiation \citep{kannan2011recurrent}. The developmental increase may reflect the establishment of new transcriptional neighborhoods as chromatin architecture reorganizes during cell fate specification \citep{clark2019single}.

The finding that 61.9\% of chimeric RNAs involve two protein-coding parental genes, with 11.2\% predicted to be in-frame, suggests that a subset of chimeric RNAs in the developing retina may encode novel fusion proteins with functional significance. However, the lack of correlation between chimeric RNA expression and parental gene expression indicates that chimeric RNA generation is not simply a byproduct of high transcriptional activity but involves regulated biogenesis mechanisms that remain to be elucidated.

Perhaps the most striking finding is the dual effect of CTCL knockdown on retinal organoid development: obstruction of neural retinal differentiation coupled with promotion of RPE fate. In the developing eye, the optic vesicle gives rise to both neural retina and RPE through a binary fate decision regulated by opposing transcription factor networks \citep{fuhrmann2014eye,bharti2006transcription}. Our results suggest that CTCL participates in this binary decision, normally favoring neural retinal fate. When CTCL is depleted, the balance shifts toward RPE, as evidenced by upregulation of MITF, RPE65, BEST1, and TYR.

The unexpected involvement of EMT and extracellular matrix pathways, rather than the anticipated Wnt signaling, in the CTCL knockdown phenotype underscores that chimeric RNAs can acquire functions independent of their parental genes \citep{wu2019chimeric}. This functional independence is consistent with the concept that chimeric transcripts represent a distinct functional class that may adopt novel protein-protein interaction capabilities or regulatory functions through the juxtaposition of domains from different genes.

The lack of Wnt pathway modulation despite CTNNBIP1 being a known Wnt inhibitor, combined with the absence of parental gene expression changes upon CTCL knockdown, strongly argues against a simple model in which the chimeric RNA merely sequesters or modulates parental gene function. Instead, CTCL appears to have evolved a distinct molecular mechanism, possibly through its in-frame fusion protein isoform, that impinges on EMT and ECM regulation in ways that influence the balance between neural and epithelial (RPE) fates.

Our study has several limitations. First, retinal organoids, while powerful, do not perfectly recapitulate in vivo retinal development, particularly lacking vascularization and RPE-photoreceptor interactions present in the native eye \citep{capowski2019reproducibility}. Second, our chimeric RNA detection relied on a single spanning unique read threshold, which provides broad sensitivity but may include some false positives. Third, the functional studies focused on a single chimeric RNA (CTCL), leaving the functions of the vast majority of detected chimeric RNAs unexplored.

\section{Conclusion}

We have established the first chimeric RNA expression atlas of human retinal organoid development, revealing dynamic regulation of chimeric RNAs across eight developmental stages from D0 to D120. The developmental increase in intra-chromosomal chimeric RNAs, the predominance of protein-coding parental gene pairs, and the functional demonstration that CTCL regulates the neural retina-RPE fate decision collectively establish chimeric RNAs as a previously unrecognized regulatory layer in retinogenesis. These findings open new avenues for understanding retinal development and may inform strategies for directing retinal organoid differentiation for therapeutic applications.

\section*{Data Availability}

RNA-seq data have been deposited in the Gene Expression Omnibus (GEO) database. Code for chimeric RNA analysis is available on GitHub.

\section*{Acknowledgments}

We thank the members of the laboratory for helpful discussions and technical assistance. This work was supported by institutional funds and research grants.

\bibliographystyle{plainnat}
\bibliography{references}

\end{document}
